\section{Introduction}

Neural activity at the mesoscopic scale---the coordinated fluctuations of local populations of neurons spanning multiple cortical
regions---plays a central role in sensory processing, motor control, and cognition~\citep{Churchland2012, Musall2019, Stringer2019}.
These dynamics unfold over spatial scales of hundreds of microns to millimeters and temporal scales of tens to hundreds of milliseconds,
an intermediate regime between the activity of single neurons and the global patterns observed in whole-brain imaging such as fMRI.
Understanding how mesoscale cortical dynamics give rise to behavior remains a fundamental challenge in systems neuroscience.

Wide-field calcium imaging has emerged as a powerful tool for monitoring mesoscale cortical activity in behaving
mice~\citep{Wekselblatt2016, Allen2017, Makino2017}. By combining genetically encoded calcium indicators such as GCaMP with chronic
cranial windows, wide-field imaging enables simultaneous recording of population-level fluorescence signals across the entire dorsal
cortex at frame rates sufficient to resolve behaviorally relevant dynamics~\citep{Ren2021, Ghanbari2019}. Studies using this approach
have revealed cortex-wide patterns of activation during sensorimotor tasks~\citep{Musall2019, Pinto2019}, the emergence of coordinated
cortical sequences during motor learning~\citep{Makino2017}, and the spatiotemporal structure of spontaneous cortical
activity~\citep{Mohajerani2013, Vanni2014}. Collectively, these observations suggest that the spatiotemporal dynamics of mesoscale
cortical activity carry rich information about the computational state of the brain.

A complementary line of work has sought to formalize neural population activity in terms of dynamical systems, in which the trajectory
of population-level activity evolves according to mathematically tractable rules. Seminal studies in motor cortex demonstrated that
neural population dynamics during reaching could be captured by low-dimensional, approximately linear
models~\citep{Churchland2012, Cunningham2014}. This dynamical systems perspective has since been applied across many brain areas and
behavioral contexts, revealing that complex single-neuron responses often reflect the projection of simpler population-level dynamics
onto individual cells~\citep{Vyas2020, Saxena2019}. Crucially, when viewed through population-level measurements, high-dimensional
nonlinear neural systems often admit locally low-dimensional or approximately linear representations---an observation that opens the
door to control-theoretic approaches for manipulating neural activity.

While observation of mesoscale dynamics has advanced rapidly, establishing causal links between these dynamics and behavioral outcomes
remains challenging. Optogenetics has transformed the study of neural circuits by enabling temporally precise, cell-type-specific
manipulation of neural activity~\citep{Deisseroth2011, Fenno2011}. Most optogenetic experiments employ open-loop stimulation, in which
predetermined patterns of light are delivered irrespective of the ongoing neural state. Although open-loop approaches have yielded
fundamental insights into circuit function, they are inherently limited in their ability to account for the variability and
state-dependence of neural responses~\citep{Grosenick2015}. Cortical activity fluctuates substantially from trial to trial due to
ongoing spontaneous dynamics, arousal, and behavioral state~\citep{Harris2011}, meaning that the same open-loop stimulus can produce
widely varying neural and behavioral effects depending on the brain state at the time of stimulation.

Closed-loop paradigms, in which stimulation is continuously adjusted based on real-time measurements of neural activity, offer a
principled framework for overcoming these limitations. From a control-theoretic perspective, feedback enables a system to compensate
for unmodeled disturbances and maintain a desired output despite variability in the underlying activity state or
dynamics~\citep{Astrom1995}. In the context of neural manipulation, closing the loop provides two key advantages: it allows
stimulation to adapt to the current brain state, improving the precision of neural modulation; and it enables the experimenter to
drive neural activity toward specific target trajectories, creating opportunities for causal tests that are not possible with
open-loop stimulation alone~\citep{OShea2018}.

Several groups have demonstrated the feasibility and utility of closed-loop optogenetic control at the level of single neurons and
local field potentials. \citet{Newman2015} introduced the ``optoclamp,'' a proportional--integral (PI) controller that locked the
firing rate of individual neurons to stationary targets in cultured networks and in the anesthetized rat thalamus.
\citet{Bolus2018} developed model-based design strategies for dynamic closed-loop optogenetic neural control in vivo, and subsequently
demonstrated state-space optimal feedback control of single-neuron firing rates in awake, head-fixed mice using linear-quadratic
optimal control~\citep{Bolus2021}. At the network level, closed-loop optogenetic stimulation has been used to modulate oscillatory
dynamics in brain slices and in non-human primates~\citep{Nicholson2018, Cagnan2023}. \citet{Clancy2014} showed that closed-loop
optogenetic feedback could train mice to modulate neural activity in motor cortex, establishing that neural circuits can be shaped
through activity-dependent stimulation. A comprehensive review by \citet{Grosenick2015} outlined the theoretical and practical
foundations for closed-loop optogenetic control, emphasizing its potential for causal circuit dissection.

Despite these advances, most prior closed-loop optogenetic studies have operated at the level of single neurons or local field
potentials, using electrophysiological recordings as the feedback signal. The application of closed-loop control to mesoscale cortical
population dynamics---using wide-field calcium imaging as the readout---remains largely unexplored. This gap is significant because
mesoscale dynamics reflect the integrated activity of large neuronal populations and are directly relevant to the kinds of distributed
cortical computations that underlie perception and behavior. Moreover, recent technical developments have made it possible to perform
simultaneous wide-field calcium imaging and spatially targeted optogenetic stimulation across the dorsal cortex of awake
mice~\citep{Matveev2024}, providing the experimental infrastructure needed for real-time, mesoscale closed-loop control.

Here, we demonstrate a closed-loop optogenetic control framework for regulating mesoscale cortical population activity in head-fixed
mice. Using wide-field calcium imaging, we extract a real-time readout of regional cortical activity ($\Delta F/F$) and use this
signal to drive a proportional--integral (PI) feedback controller that modulates targeted optogenetic stimulation of a neighboring
cortical region. We first characterize the input--output relationship between optogenetic stimulation and trial-averaged cortical
responses, showing that this relationship is well approximated by a linear mapping---a finding that enables principled, model-based
controller design. We then implement a PI controller with feedforward compensation and demonstrate that closed-loop stimulation
improves reference tracking accuracy and reduces trial-to-trial variability compared to open-loop stimulation across multiple
experimental sessions and animals. The reduction in variability is interpreted as an active disturbance rejection phenomenon,
wherein feedback suppresses the influence of spontaneous cortical fluctuations on the driven response. Collectively, these results
establish that stationary, approximately linear trial-averaged dynamics can be exploited for modeling and controlling regional
cortical activity, and demonstrate the potential of feedback-based perturbation paradigms for causally investigating cortical
population dynamics.
